%meeting 3  26 - 08 - 26
Kindly clean this up, this is a relativity lecture. Add proper equations, fix spelling, flesh out the arguments, and know what shit means, as best you can. Make \subsubsections in appropriate points.
Keept it mostly verbatim in the clear sections.
Use the physics package macros when you can. Return in latex and your own ddisplay. 

\subsection{Relativity without Light}

This is in contrast to standard methods of deriving SR.

Suppose framce of references S and S', where S is fixed and S' is moving at a speed $v$
graph of space.

In galilean 

Then
\begin{align*}
	x' = \bar[X](x,t;v)
	t' = \bar{T}(x,t;v)
\end{align*}

Task: Determine function $\bar X$ and $\bar T$

Goal - the transpfolation lorentz 
x' 
t' 
gamma 
beta
(x' t) = gramma ( 1 -v \\ -beta 1 ) (x,t)


Assujme relativity principle
S and S' are equivalent

X bar and T bar are transformation functions. The role is to find them because they are undetermined functions and we want to find them out.

x = \bar{X} (X, t', -v)
t = \bar T (x',t', -v)

This also assumes the isotropgy od space

Hence (eq 1)

\begin{align}
	x = \bar X ( \bar X(t,x,v), \bar{T}(x,t,v),-v )
	t - \bar T ( bar X(t,x,v), \bar{T}(x,t,v), - v )
\end{align}

I sotrophy of space means no difference between +x and -x directions.

Reversing the t direction amounts to changing 
x\to-x 
v\to - v

x' = -\bar X ( -x, t, -v )
t' = \bar T (-x, t, -v)

\bar X (-x,t,-v) = -\bar X (x,t,v)
\bar T (-x,t,-v) = -\bar T (x,t,v)

The invariance of straight lines

Free particles must move along straight lines, in both $S$ amd S'
The first law of motion of Newnton, principle of inertia

let $x^\mu = (t,x)$ amd $x'^\mu = (t',x')$  

Suppose $x^\mu = (t(\tau), x(\tau))$ and $x'\mu = (t'(\tau),t'(\tau))$

How does oen express straight lines in S
\begin{equation}
	\dv[2{x^\mu}{}\tau^2 = 0
\end{equation}
we impose that 
\begin{equation}
	\dv[2{x^\mu}{}\tau^2 = 0 \implies \dv[2]{x'^\mu}{\tau^2} 
\end{equation}


We can then write
\begin{equation}
    \dv{}\tau{}x'^\mu = \dv{\bar X^\mu}{x^\nu} \dv{x^\nu}{\tau}   
\end{equation}

HEre, we impose the einstain notation convention. 

\begin{equation}
	\dv[2]{}\tau{}x'^\mu = \dv{}{\tau}( \dv{\bar X^\mu}{x^nu} ) \dv{x^\nu}{\tau} + \pdv{\bar X^\mu}{x^\nu} \pdv{x^\nu}{\tau}   
\end{equation}

We demand that the second term vanishes to zero. However, since the LHS = 0 , then we find that 
\begin{equation}
\dv{}{\tau}( \dv{\bar X^\mu}{x^nu} ) \dv{x^\nu}{\tau}
\end{equation}
Then
\begin{equation}
	0 = \pqty{
		\pdv[2]{\bar x^\mu}{\partial x^\nu \partial x^\sigma} \dv{x^\nu}{\tau} \dv{x^\sigma}{\tau}   
	}
\end{equation}

Hence
\begin{equation}
    \frac{\partial \bar X^\mu}{\partial x^\nu \partial x^\sigma} = 0 
\end{equation}

We only get that if we assume $\bar X$ is a linear function of $x^\mu $ or (t,x).

$\bar X (t,x,v) = A(v)x+ B(v)t+ E(v)$ 
$\bar T (t,x,v) = C(v)x+ D(v)t+ E(v)$ 

Weve not done anything but write a straight line in general.

There is no preferred zero for x in t in both S' and S, that's the homogeneity of space time.

Assumption: Space and time are homogeneous.

We can choose for the origins of $S$ and $S'$ 's origins correspand

We demand that the origin of S gets mapped to the origin of S'
\begin{equation}
    x = t = 0 \iff x' = t' = 0
\end{equation}

This is not really essential, but it makes calculations easier, making E  and F determinable (they would be 
difficult to obtain) so we can simplify E and F to be 0.

Now we only require four free functions

$\bar X (t,x,v) = A(v)x+ B(v)t
$\bar T (t,x,v) = C(v)x+ D(v)t
The problem now is finding four functions of $V$. 

\[
	-A(-v) x + B(-v) t = -A(v)x - B(v)t\\
	-C(-v) x + D(-v) t = C(v)x + D(v) t
\]

recall from the condition that 
\bar X (-x,t,-v) = -\bar X (x,t,v)
\bar T (-x,t,-v) = -\bar T (x,t,v)

This mist be true for all of X and T. THe only way this can hold true if the eq is equal. These holds four impositions

Therefore
\[
	A(-v)  = A(v)
\]
\[
	D(-v)  = D(v)
\]
\[
	B(-v) = -B(v)
\]
\[
	C(-v) = - C(v)
\]

Therefore a and D need to be even functions ,and B and C need to be odd functions.

we can then write equation 1 as
\[
	x = A(-v)\bar X(t,x,v) + B(-v) \bar T(t,x,v)\\
	t = C(-v)\var X(t,x,v) + D(0v) \bar T(x,t,v)
\]


There is still a \bar X and \bar T that appears. We can write that moer out

\[
	x = A(-v) ( A(v)x + B(v)t + B(-v)) (C(v)x + D(v) t)
	  + D(-v)(C(v)x + D(v)t)
	t = C(0v)( A(v) x + B(v)t ) + D(-v)(C(v)x + D(v))
\]

We impose the odd and even conditions on the eqn. We can then noew write
\[
	x = (A^2 - BC) x _ (AB - BD) t\\
	t = (-AC+CD)x + (-BC + D^2)
\]
We use the argument that x and t are arbitrary. We can then find
\[
	A^2 - BC = 1\\
	B(A-D) = 0\\
	C(A0D) = 0\\
	D^2 -BC = 1
\]
So far, no magic. We just used natural assumptions and principle of relativity. 
We focus on the two conditions that equat to zero. 
We have two options
A \neq D or A = D

If A\neq D \implies B = C = 0, we are forced to conclude.

if that is true,then
A=\pm 1
D =\mp 1

What happesn to the transformation?

x' = \pm x 
t' = \mp t

What does that not mean? It's a trivial transformation, but we find that S and S' has these posibilities.
It's a cousin of the identity transformation, it satisfies the conditions that we constrain, but this is tribvial, and not interesting

But when A = D, B and C no longer has to be equal

Let's look at the other case

if A = D

Notice that these two collapse to the same condition

A = D \implies C = A^2 - 1  / B

Now we have three constraints in essence. Now, we only have two free functions. since onece you know A and B you know C

A and B are the only free functions left 

Note that (Consider the trajectory of S' and S) S' moves in S in a fixed velicity v. So its trajectory in S is x = vt
Again, that is for the spatial origin. That means that entire line needs to be mapped to x' = 0.

Now, we impose the C = A^2 -1 / B

x' = Ax + Bt
t' - (A^2 - 1 / B) + At

hence
0 = Avt + Bt \\
B = -vA

how we only have one free functions

Thus we have redcued it to the following

\bqty{
	x'\\t'
} =
\bqty{
	A & -vA \\
	- \frac{(A^2-1)}{vA} & A
}
\bqty{x\\t}

Now we only need to know A, except we dont know what properties of A is

That matrix needs to reduce to the identiy matrix when v = 0

Thus we impose

A(v)\to 1, v\to 0. Since fraction, it must vanish sufficiently fasst.
therefore, A^2 - 1 \approx \order{v^2}

What happens if there's a third observer?

Consider a new frame S'' moving with a speed u wrt S' 

It's the same framework. Now we relate S'' to S.

So now we have

\bqty{
	x'\\t'
} 
= 
\bqty{
	A(u) & -uA(u) \\
	- \frac{(A(u)^2-1)}{uA(u)} & A(u)
} 
\bqty{
	A & -vA \\
	- \frac{(A^2-1)}{vA} & A
}
\bqty{x\\t}

Now, we demand closure under composition. That must be of the samme form as before. 
Therefoer it must be the same matrix form but with a different argument. Say that S'' travels at a speed w wrt S

Now, we multiply
\bqty{
	A(u) A(v) + (A^2(v) - 1) \frac{u A(u)}{v A(v)}  & - (u+v)A(u) A(v)\\
	-(A^2(u)-1) \frac{A(v)}{uA(u)} - (A^2(v)-1 ) \frac{A(v)}{vA(v)} & A(u) A(v) + (A(u)^2 - 1 )  \frac{v A(v)}{u A u} \\
}

From the form, we require that M^{00} and M^{11} are equal.

Therefore we find
\begin{equation}
    A^2(v) - 1 / v^2 A^2(u) = A^2(u) - 1 / u^2 A^2 u
\end{equation}

Notice that the lhs depends entirely on v and the rhs depends entirely on u. As such, the whole combination needs to be constant. No matter
what u or v you put, it has to equal. hence

A^2  - 1 / v^2 A^2 = k

A = \frac{1}{\sqrt(1-kv^2)} 

k has to have a unit of 1/v^2. 

the M&{01} term becoems A^2 -1 / va = -kA V

Hecne

\bqty{
	x'\\t'
}
=
\bqty{
	A & - vA\\
	-AkV & A
}
\bqty{
	x\\t
}
\bqty{
	x\\t
}

Factor out A

\bqty{
	x'\\t'
}
=
\frac{1}{\sqrt{1-kv^2}} 
\bqty{
	1 & - v\\
	-kV & 1
}
\bqty{
	x\\t
}
\bqty{
	x\\t
}
this is exactly the form of the lorentz transfomration.

Thus we can find out that 
\begin{equation}
    W = \frac{u+v}{1 + kuv} 
\end{equation}
this is the relativistic velocity addition formula.

If V = 1/sqrt{k}, the $w= \frac{1}{\sqrt ko} $

Therefore, ralativity demands that some speed to be invariant.
There must exist some invariant speed. That speed needs to be v = 1/sqrt k

We dont need light. It does not need anything.
This is the biggest principle to come ouyt. Relativity works from pure logic alone.
All physics must bow down to this principle

v tends to infinity id k = 0. therefore

we arrive to the galilean transfomrations.

It does not need to be the speed of light, it could be tghe speed of anything, ut just happens to be light.
It does not mean that relativity demands Em to be true for it to be true. 

Empirically, we find em has an invariant speed. Em is a calibrator for relativity.

You can imagine a consistent universe that the invariant speed is infinity. 

Carrolian Geometry.  Causality 

For anything finite, reality tends to be a light cone. 
The light cone opens up.
at V = 0, the light cone closes, and you cannot move

Infinity, finite, or zero. 
Galileo, Lorentz, Carol

Thus, we prove that relavity, by logic alone, demans that an invariant speed must exist in the universe.
